\section{Equilibrium physical windows and relative conditional laws}
\label{sec:stationary-physical-conditioning}
A typical physical observation does not start at a section collision.
The section state seen from equilibrium is biased by the duration
of its return excursion. We retain this bias exactly. A wider
collision-clock Fourier band supplies the denominator, while the
return-clock moment estimates compare the same stationary physical
event with its completed-return descriptions. This avoids assuming
bounded variation of a length-biased entrance distribution.

\subsection{Two descriptions of one equilibrium probability}
Write $\Theta_R(y)=\tau_R^*(y)$, $r_R^*(y)=r(y)$, and
$\bar\Theta_R=\int\Theta_R\dd\nu_R^*=\bar\tau_R/c_*$. By
Proposition~\ref{prop:induced-suspension}, equilibrium flow is represented on
\begin{equation}\label{eq:stationary-return-suspension}
 \mathcal X_R=\{(y,s):y\in Y_R^*,\ 0\le s<\Theta_R(y)\},\qquad
 \dd\mathbb P_R(y,s)=\bar\Theta_R^{-1}\dd\nu_R^*(y)\dd s.
\end{equation}
The physical observation starts at the flow state reached after time
$s$ from $y$. For $0\le j\le r(y)$ put
$\tau_{j,R}^{\rm sum}(y)=\sum_{i<j}\tau_R(T_R^iy)$.
There is a unique $0\le j<r(y)$ such that
$\tau_{j,R}^{\rm sum}\le s<\tau_{j+1,R}^{\rm sum}$, apart from
null endpoint conventions. Define the current outgoing collision and
its elapsed flight age by
\[
 x=T_R^jy,\qquad u=s-\tau_{j,R}^{\rm sum}(y).
\]

\begin{lemma}[The two length biases and exceptional-set transfer]
\label{lem:stationary-length-bias}
Under $\mathbb P_R$, the marginal law of $y$ is
$(\Theta_R/\bar\Theta_R)\nu_R^*$, whereas the law of $(x,u)$ is
\begin{equation}\label{eq:collision-suspension-marginal}
 \bar\tau_R^{-1}\one_{\{0\le u<\tau_R(x)\}}\dd\nu(x)\dd u.
\end{equation}
In particular the outgoing collision marginal is
$\sigma_R\nu$, where $\sigma_R=\tau_R/\bar\tau_R$ is a probability
density with uniformly bounded supremum and BV norms.
There are uniform $C,c>0$ such that, for measurable $E\subset Y_R^*$
and $L\ge0$,
\begin{equation}\label{eq:stationary-bias-bounds}
 \mathbb P_R\{y\in E\}\le C\nu_R^*(E)^{1/2},\qquad
 \mathbb P_R\{\Theta_R(y)>L\}\le Ce^{-cL}.
\end{equation}
\end{lemma}
\begin{proof}
Integrating $s$ in \eqref{eq:stationary-return-suspension} gives the
first marginal. For an integrable test $F$ of $(x,u)$, partition each
excursion into its actual collision flights. The tower identity gives
\begin{align*}
 \int F(x,u)\dd\mathbb P_R
 &=\frac1{\bar\Theta_R}\int_{Y_R^*}\sum_{j<r(y)}
       \int_0^{\tau_R(T_R^jy)}F(T_R^jy,u)\dd u\dd\nu_R^*(y)\\
 &=\frac1{c_*\bar\Theta_R}\int_M\int_0^{\tau_R(x)}F(x,u)\dd u\dd\nu(x).
\end{align*}
Since $c_*\bar\Theta_R=\bar\tau_R$, this proves the second marginal.
Lemma~\ref{lem:physical-bv} and the positive lower bound for
$\bar\tau_R$ give the density bounds. No assertion that
$\Theta_R/\bar\Theta_R$ is bounded-BV is needed or made.

The return roof has a uniform exponential moment by
Theorem~\ref{thm:exponential-return}. In particular
$\|\Theta_R\|_{L^2(\nu_R^*)}$ is uniform. Cauchy--Schwarz proves
the first estimate in \eqref{eq:stationary-bias-bounds}. Integrating
the exponential roof tail after multiplication by $\Theta_R$, and
reducing the tail exponent, proves the second. The positive normalizing
mean is bounded away from zero. This treats the physical-time length
bias, not merely the collision-count bias by $r(y)$.
\end{proof}

\subsection{Coupling collision, return and physical observations}
Let $V_R(t)=(W_R(t),C_R(t)-t/\bar\tau_R)$ be the actual physical
increment during the interval of length $t$ starting at $(y,s)$.
The lattice convention is the one used throughout the article; any
bounded cell-position term is included in the following error bounds.
Retain $B_R$ from \eqref{eq:section-start-physical-vector}. Set
\begin{align}
 m_R(t)&=\lfloor t/\bar\tau_R\rfloor,&
 n_R(t)&=\lfloor c_*t/\bar\tau_R\rfloor,\notag\\
 L_R(t;y,s)&=\max\{l\ge0:T_{l,R}(y)\le t+s\},\notag\\
 Z_t^{\rm coll}&=t^{-1/2}B_R\sum_{i<m_R(t)}h_R(T_R^ix),&
 Z_t^{\rm ret}&=t^{-1/2}B_RU_{n_R(t),R}(y),\notag\\
 Z_t^{\rm last}&=t^{-1/2}B_RU_{L_R(t;y,s),R}(y),&
 Z_t^{\rm phys}&=t^{-1/2}V_R(t).
 \label{eq:stationary-four-observations}
\end{align}
These are random variables on the same probability space
$(\mathcal X_R,\mathbb P_R)$. The return reference starts at the
actual section origin $y$ preceding the stationary observation, not
at a newly sampled independent section state.

\begin{proposition}[Uniform stationary observation coupling]
\label{prop:stationary-observation-coupling}
For each fixed $P>0$ there is $C_P$ such that, for all large $t$,
\begin{equation}\label{eq:stationary-observation-coupling}
 \mathbb P_R\left\{\max_{j\in\{{\rm ret},{\rm last},{\rm phys}\}}
 |Z_t^j-Z_t^{\rm coll}|>C_Pt^{-1/8}\right\}\le C_Pt^{-P}
 \qquad(R\in I).
\end{equation}
\end{proposition}
\begin{proof}
First compare the collision reference with the physical vector.
The age $u$ is bounded by the uniform one-flight horizon. If
$\ell(t)$ is the completed-flight index measured from $x$, then
its threshold differs from $t$ by at most that horizon. The flight
roof is bounded-BV and has mean $\bar\tau_R$. Fixed collision
moments and the monotone inverse-clock equivalences therefore give,
with $m=m_R(t)$, $q=\lceil t^{5/8}\rceil$, and each fixed $p$,
\[
 \mathbb P_R\{|\ell(t)-m|>q\}\le C_p t^{p/2}q^{-p}
                                            \le C_pt^{-p/8}.
\]
Indeed the outgoing collision density is bounded by
Lemma~\ref{lem:stationary-length-bias}, so stationary collision
moment estimates apply by domination; independence of $u$ and $x$
is not used. Off this event the centered collision-sum difference
lies in the deterministic length-$q$ windows adjacent to $m$.
Lemma~\ref{lem:fixed-even-maximal}, at a sufficiently large fixed
even order, bounds the probability that its size exceeds $t^{3/8}$
by $C_p q^{p/2}t^{-3p/8}\le C_pt^{-p/16}$.
The initial and final partial flights contribute $O(1)$ after applying
$B_R$, because $B_Rh_R=(\kappa_R,1-\tau_R/\bar\tau_R)$.
Choose $p>16P$. This proves the physical-to-collision comparison.

For the return reference, apply
Proposition~\ref{prop:physical-window-coupling} under $\nu_R^*$,
with its exceptional probability exponent chosen larger than $2P$.
The first estimate in \eqref{eq:stationary-bias-bounds} transfers
this exceptional set to $\mathbb P_R$ with probability $O(t^{-P})$.
The actual interval from time $s$ to $s+t$ differs from the
section-start interval from zero to $t$ by two end intervals of
length at most $s\le\Theta_R(y)$. The positive minimum free-flight
length and bounded one-flight displacement give, deterministically,
\[
 |V_R(t)-V_R^Y(t;y)|\le C(1+\Theta_R(y)).
\]
This includes collision-count centering and cell-position conventions.
The exponential tail in \eqref{eq:stationary-bias-bounds} makes this
error at most $Ct^{3/8}$ outside a smaller exceptional set. Thus the
return reference and the physical vector obey the required comparison.

For the last-return reference, its difference from the physical vector
is bounded before normalization by
\[
 C\bigl(1+\Theta_R(y)+r_R^*((F_R^*)^{L_R(t;y,s)}y)\bigr).
\]
The initial partial return contributes the roof term and the final
partial return contributes the last term. On $\Theta_R\le t$,
the positive minimum flight length gives $L_R(t;y,s)\le C(1+t)$.
Under $\nu_R^*$ the union of the events
$r_R^*\circ(F_R^*)^j>t^{3/8}$ over these deterministic indices has
probability at most $C(1+t)e^{-ct^{3/8}}$, by return invariance.
The first bound in \eqref{eq:stationary-bias-bounds} transfers this
union to $\mathbb P_R$; the event $\Theta_R>t$ has exponential
probability as well. This proves the last-return comparison.
A triangle inequality and union bound now give
\eqref{eq:stationary-observation-coupling} for all three variables
against the common collision reference. All moment orders are fixed
for a chosen $P$.
\end{proof}

\subsection{The stationary denominator and relative replacement}
Let $h=t^{-2/25}$ and $B_t=B_3(\xi,(h,h,h))$, $|\xi|\le A$.
For each $j\in\{{\rm coll},{\rm ret},{\rm last},{\rm phys}\}$ define
\begin{equation}\label{eq:stationary-four-events}
 A_t^j=\{Z_t^j\in B_t\}.
\end{equation}
The physical event is an observation at the original deterministic
time $t$, with exact integer displacement and collision-count
intervals of physical half-width $t^{21/50}$.

\begin{theorem}[Stationary physical windows and relative conditional laws]
\label{thm:stationary-physical-windows}
Uniformly for $R\in I$, $|\xi|\le A$, and large $t$, each event in
\eqref{eq:stationary-four-events} satisfies
\begin{equation}\label{eq:stationary-window-denominator}
 \mathbb P_R(A_t^j)=8t^{-6/25}g_{\mathcal V_R}(\xi)
            \left(1+O_A(t^{-23/2800}\sqrt{\log(2+t)})\right).
\end{equation}
For $j={\rm coll},{\rm last},{\rm phys}$,
\begin{equation}\label{eq:stationary-relative-event-bound}
 \frac{\mathbb P_R(A_t^j\mathbin\triangle A_t^{\rm ret})}
                 {\mathbb P_R(A_t^{\rm ret})}
 \le C_A t^{-23/2800}\sqrt{\log(2+t)}.
\end{equation}
The conditional laws $\mathbb P_R(\,\cdot\mid A_t^j)$ and
$\mathbb P_R(\,\cdot\mid A_t^{\rm ret})$ on the full stationary
initial-state space have total variation distance bounded by the
same right side with a changed constant. Thus the comparison holds
for every common bounded measurable statistic of the entire orbit.
\end{theorem}
\begin{proof}
Take $m=m_R(t)$ and the initial density
$\sigma_R=\tau_R/\bar\tau_R$. Apply
Theorem~\ref{thm:projected-collision-windows} with
$P_{m,R}=\sqrt{m/t}\,B_R$, augmenting by the fourth row $(0,0,0,1)$.
The ratios $m/t$, $\bar\tau_R$ and their reciprocals have uniform
bounds for large $t$, so the augmenting matrix and its inverse do
also. The covariance is
\[
 (m/t)B_R\Gamma_RB_R^{\mathsf T}
                         =\mathcal V_R+O(t^{-1}).
\]
The window half-width $t^{-2/25}$ is comparable to $m^{-2/25}$.
Uniform ellipticity controls the change in the Gaussian density from
this $O(t^{-1})$ perturbation. Replacing its integral on the window
by $8h^3$ times its central value costs relative $O_A(h)$.
This proves \eqref{eq:stationary-window-denominator} for
$A_t^{\rm coll}$ and gives a lower bound $c_Ah^3$.

Choose $P=2$ in
Proposition~\ref{prop:stationary-observation-coupling} and put
$d_t=C_2t^{-1/8}$. If $A_t^j$ and $A_t^{\rm coll}$ differ outside
the exceptional set, then $Z_t^{\rm coll}$ belongs to the shell
between boxes of half-widths $h+d_t$ and $h-d_t$ about $\xi$.
For large $t$ those widths are between $h/2$ and $2h$.
The projected collision-window theorem applies to both boxes with
the same constants. Subtract their probabilities before taking
absolute values. The Gaussian shell mass is at most $C_Ah^2d_t$;
the total envelope comparison error is at most $C_Ah^3\varepsilon_t$.
Consequently
\begin{equation}\label{eq:stationary-shell-budget}
 \frac{\mathbb P_R(A_t^j\mathbin\triangle A_t^{\rm coll})}
                         {\mathbb P_R(A_t^{\rm coll})}
 \le C_A\left(t^{-9/200}+\varepsilon_t+t^{-2+6/25}\right)
 \le C_A\varepsilon_t.
\end{equation}
This also proves the denominator asymptotic for each $j$.
The symmetric-difference triangle inequality and the resulting
comparison of the denominators give
\eqref{eq:stationary-relative-event-bound}. No Fourier estimate of
an unresolved thin boundary slab is assumed.

For measurable events $A,B$ with $p=\mathbb P_R(A)>0$ and
$\delta=\mathbb P_R(A\triangle B)\le p/2$, direct subtraction of
$\one_A/p$ and $\one_B/\mathbb P_R(B)$ gives
$d_{\rm TV}(\mathbb P_R(\cdot\mid A),\mathbb P_R(\cdot\mid B))
\le2\delta/p$. Apply this to the proved event comparison. The
expectation difference for a common bounded statistic is at most
twice its supremum norm times that distance. This is a comparison
of conditional measures, not a claim of a conditional invariance
principle or of independence from the event.
\end{proof}

\begin{corollary}[Additional bounded path selection]
\label{cor:stationary-weighted-selection}
Let $w_t\ge0$ be any measurable function on $\mathcal X_R$, possibly
depending on the whole path, with $\|w_t\|_\infty\le B$ and
\[
 \int_{A_t^{\rm ret}}w_t\dd\mathbb P_R
                   \ge a_t\mathbb P_R(A_t^{\rm ret}),\qquad a_t>0.
\]
If $B\varepsilon_t/a_t\to0$, the conditional laws with densities
proportional to $w_t\one_{A_t^j}$ and $w_t\one_{A_t^{\rm ret}}$
have total variation distance $O_A(B\varepsilon_t/a_t)$.
\end{corollary}
\begin{proof}
Their unnormalized difference has $L^1$ norm at most
$B\mathbb P_R(A_t^j\triangle A_t^{\rm ret})$.
Divide by the stated lower bound for the weighted return denominator,
then use the same normalization inequality as in the preceding proof.
The lower bound involving $a_t$ is an explicit hypothesis, not a
weighted raw denominator theorem inferred from the unweighted one.
\end{proof}

\paragraph{Relation to the original raw problem.}
Theorems~\ref{thm:wide-collision-band} and
\ref{thm:stationary-physical-windows} supply a collision-clock
Fourier estimate and a relative physical-event theorem under the
actual equilibrium law. They are not statements under the
unbiased section-start law. The physical half-width $t^{21/50}$
is smaller than the section-start half-width $t^{23/50}$ above,
but still grows: it is not a singleton lattice label. The probability
$t^{-6/25}$ and the relative error in
\eqref{eq:stationary-relative-event-bound} refer to these exact
windows and to no substituted event. The bounded path-selection
corollary states its own additional denominator requirement.

The full prescribed-return-count complement, the pointwise common
signed correction, the long-time inverse-coarea derivative budget,
and the original microscopic conditioning theorem remain distinct
requirements of the same raw mixed-density problem. In particular
removing the collision stopping loss has not removed the return-count
stopping loss from Theorems P and Q. The common raw correction is still
bounded only in the signed window norm proved in
Theorem~\ref{thm:averaged-common-raw-remainder}, unless its stronger
pointwise norm is established separately.
